\documentclass[11pt,spanish]{article}

% =========================================================
% PLANTILLA BASE PARA INFORMES ACADÉMICOS / TÉCNICOS
% =========================================================

% --- Idioma y codificación ---
\usepackage[spanish,es-tabla]{babel}
\usepackage[T1]{fontenc}
\usepackage{textcomp}
\usepackage{newunicodechar}
\newunicodechar{₡}{\textcolonmonetary}

% --- Márgenes / papel ---
\usepackage[
    left=2.5cm,
    right=2.5cm,
    top=2.5cm,
    bottom=2.5cm,
    headsep=0.5cm,
    footskip=0.8cm
]{geometry}

% --- Matemáticas ---
\usepackage{amsmath}

% --- Tablas, enlaces y elementos visuales ---
\usepackage{array}
\usepackage{tabularx}
\usepackage{booktabs}
\usepackage{longtable}
\usepackage{graphicx}
\usepackage{float}
\usepackage{xcolor}
\usepackage{tikz}
\usetikzlibrary{positioning,shapes.geometric,arrows.meta,calc}
\usepackage{enumitem}
\usepackage{ragged2e}
\usepackage[hidelinks]{hyperref}
\usepackage{adjustbox}

% --- Encabezado y pie de página ---
\usepackage{fancyhdr}
\usepackage{lastpage}

% --- Interlineado ---
\linespread{1.2}
\sloppy

% =========================================================
% DATOS DEL DOCUMENTO
% =========================================================
\newcommand{\Institucion}{Instituto Tecnológico de Costa Rica}
\newcommand{\Campus}{Campus Tecnológico Central Cartago}
\newcommand{\DocumentoTitulo}{Cronograma del Proyecto}
\newcommand{\Curso}{Administración de proyectos}
\newcommand{\Profesor}{Alicia Marcela Salazar Hernández}
\newcommand{\FechaDocumento}{22 de junio de 2026}
\newcommand{\PeriodoAcademico}{I Semestre}
\newcommand{\AnioDocumento}{2026}
\newcommand{\AutoresPortada}{%
    \begin{tabular}{l@{\hspace{1.5cm}}r}
        Mauricio González Prendas & 2024143009 \\
        Isaac Jiménez Blanco & 2024202804 \\
        Tamara Robles Camacho & 2024099342 \\
    \end{tabular}%
}
\newcommand{\DocHeaderLeft}{}
\newcommand{\DocHeaderRight}{Cronograma del Proyecto}
\newcommand{\DocFooterRight}{Página~\thepage\ de~\pageref{LastPage}}

% Metadatos PDF.
\title{\DocumentoTitulo}
\author{\AutoresPortada}
\date{\FechaDocumento}

% =========================================================
% CONFIGURACIÓN DE TABLAS Y UTILIDADES
% =========================================================
\newcolumntype{Y}{>{\RaggedRight\arraybackslash}X}
\newcolumntype{L}[1]{>{\RaggedRight\arraybackslash}p{#1}}
\renewcommand{\arraystretch}{1.22}
\setlength{\LTpre}{0.5em}
\setlength{\LTpost}{0.5em}

% Imagen segura: si el archivo no existe, muestra un recuadro de marcador.
\newcommand{\SafeImage}[2][0.92\textwidth]{%
    \IfFileExists{#2}{%
        \includegraphics[width=#1]{#2}%
    }{%
        \fbox{%
            \begin{minipage}[c][4cm][c]{#1}
            \centering
            Imagen pendiente:\\
            \texttt{#2}
            \end{minipage}%
        }%
    }%
}

% Figura reutilizable.
\newcommand{\EvidenceFigure}[3]{%
    \begin{figure}[H]
    \centering
    \SafeImage{#1}
    \caption{#2}
    \label{#3}
    \end{figure}%
}

% =========================================================
% ENCABEZADO Y PIE
% =========================================================
\pagestyle{fancy}
\fancyhf{}
\fancyhead[L]{\DocHeaderLeft}
\fancyhead[R]{\DocHeaderRight}
\renewcommand{\headrulewidth}{0.4pt}
\renewcommand{\footrulewidth}{0.4pt}
\fancyfoot[R]{\DocFooterRight}

% Estilo equivalente para páginas horizontales.
\fancypagestyle{widefancy}{%
    \fancyhf{}%
    \fancyhead[L]{\DocHeaderLeft}%
    \fancyhead[R]{\DocHeaderRight}%
    \renewcommand{\headrulewidth}{0.4pt}%
    \renewcommand{\footrulewidth}{0.4pt}%
    \fancyfoot[R]{\DocFooterRight}%
}

% =========================================================
% PÁGINAS HORIZONTALES PARA TABLAS ANCHAS
% =========================================================
\makeatletter
\newcommand{\SyncPageBuilderDimensions}{%
    \global\vsize=\textheight
    \global\@colroom=\textheight
    \global\@colht=\textheight
}
\makeatother

\newenvironment{widepage}{%
    \clearpage
    \hypersetup{pageanchor=false}%
    \paperwidth=297mm
    \paperheight=210mm
    \pdfpagewidth=297mm
    \pdfpageheight=210mm
    \newgeometry{left=2.5cm,right=2.5cm,top=2.5cm,bottom=2.5cm,headsep=0.5cm,footskip=0.8cm}%
    \setlength{\textwidth}{243mm}%
    \setlength{\textheight}{156mm}%
    \setlength{\linewidth}{\textwidth}%
    \setlength{\columnwidth}{\textwidth}%
    \hsize=\textwidth
    \setlength{\headwidth}{\textwidth}%
    \SyncPageBuilderDimensions
    \pagestyle{widefancy}%
}{%
    \clearpage
    \paperwidth=210mm
    \paperheight=297mm
    \pdfpagewidth=210mm
    \pdfpageheight=297mm
    \restoregeometry
    \SyncPageBuilderDimensions
    \hypersetup{pageanchor=true}%
    \pagestyle{fancy}%
}

% =========================================================
% PORTADA
% =========================================================
\newcommand{\MakeCover}{%
\begin{titlepage}
\thispagestyle{empty}
\centering

{\bfseries \huge \Institucion \par}
\vspace{0.25cm}
{\LARGE \Campus \par}

\vfill

{\huge \DocumentoTitulo \par}

\vfill

{\bfseries \LARGE Profesora \par}
{\Large \Profesor \par}

\vfill

{\bfseries \LARGE Estudiantes \par}
{\Large \AutoresPortada \par}

\vfill

{\bfseries\LARGE Fecha \par}
{\Large \FechaDocumento \par}

\vfill

{\LARGE \PeriodoAcademico \par}
{\LARGE \AnioDocumento \par}

\end{titlepage}%
}

% =========================================================
% DOCUMENTO
% =========================================================
\begin{document}
\hypersetup{pageanchor=false}

% =========================
% Portada
% =========================
\MakeCover

% =========================
% Índice
% =========================
\pagenumbering{roman}
\pagestyle{empty}
\addtocontents{toc}{\protect\thispagestyle{empty}}
\tableofcontents
\clearpage
\pagenumbering{arabic}
\hypersetup{pageanchor=true}
\pagestyle{fancy}

% =========================
% Contenido
% =========================
\section{Información general del documento}

\begin{table}[H]
\centering
\begin{tabularx}{\textwidth}{p{5cm} Y}
\toprule
\textbf{Nombre del proyecto} & Plataforma Universitaria Integrada \\
\textbf{Organización} & Universidad Tecnológica La Mejor \\
\textbf{Documento} & \DocumentoTitulo \\
\textbf{Directora del proyecto} & Tamara Robles Camacho \\
\textbf{Herramienta utilizada} & Microsoft Project \\
\textbf{Fecha de inicio planificada} & 8 de junio de 2026 \\
\textbf{Fecha final planificada} & 9 de setiembre de 2026 \\
\textbf{Presupuesto total} & ₡17,625,200 \\
\textbf{Fecha del documento} & \FechaDocumento \\
\bottomrule
\end{tabularx}
\end{table}

\section{Introducción}

El presente documento corresponde al cronograma del proyecto Plataforma Universitaria Integrada, elaborado como parte de la planificación y control de la iniciativa desarrollada para la Universidad Tecnológica La Mejor. Este cronograma permite organizar de manera estructurada las actividades necesarias para completar los entregables de gestión, el prototipo de interfaz web navegable, las acciones de supervisión y control, y las actividades finales de cierre.

El cronograma fue elaborado utilizando Microsoft Project, debido a que esta herramienta permite registrar tareas, subtareas, dependencias, duraciones, recursos, hitos, costos, responsables y ruta crítica. Por lo tanto, el documento no se limita a presentar una lista de fechas, sino que funciona como una herramienta de gestión para controlar el avance del proyecto, identificar tareas críticas y mantener coherencia entre alcance, tiempo, costo y recursos.

El proyecto se planifica con inicio el 8 de junio de 2026 y cierre el 9 de setiembre de 2026, con una duración aproximada de tres meses. El presupuesto total del proyecto se mantiene en ₡17,625,200, distribuido entre las fases de inicio, planificación, ejecución, supervisión y control, cierre, y rubros adicionales de apoyo y contingencia.

\section{Propósito del cronograma}

El propósito del cronograma es establecer la secuencia de trabajo que debe seguir el equipo para completar el proyecto dentro del tiempo y presupuesto planificados. Además, permite visualizar qué actividades se ejecutan en paralelo, qué tareas dependen de otras, quiénes son los responsables asignados y cuáles actividades pueden afectar directamente la fecha final del proyecto.

Este entregable se orienta al cumplimiento de los siguientes objetivos específicos.

\begin{itemize}[leftmargin=*, label={\textbullet}]
    \item Definir las tareas y subtareas necesarias para completar el proyecto.
    \item Establecer duraciones, fechas de inicio y fechas de finalización.
    \item Registrar dependencias y predecesoras entre actividades.
    \item Asignar recursos humanos y recursos de costo.
    \item Identificar responsables por actividad.
    \item Marcar hitos relevantes para el control del avance.
    \item Determinar la ruta crítica del proyecto.
    \item Integrar costos por actividad y por fase.
    \item Brindar una base para el monitoreo y control del proyecto.
\end{itemize}

\section{Verificación de cumplimiento del Entregable 6}

El cronograma cumple con los elementos solicitados para el Entregable 6, ya que fue elaborado en Microsoft Project y contiene los componentes principales requeridos para una planificación formal del proyecto.

\begin{table}[H]
\centering
\caption{Verificación de requisitos del Entregable 6}
\begin{tabularx}{\textwidth}{p{4.3cm} Y}
\toprule
\textbf{Elemento solicitado} & \textbf{Evidencia dentro del cronograma} \\
\midrule
Tareas & Se muestran en la tabla del Diagrama de Gantt, donde cada fila representa una actividad del proyecto. \\
Subtareas & Se reflejan mediante actividades agrupadas dentro de fases como Inicio, Planificación, Ejecución, Supervisión y control, Cierre y Rubros adicionales. \\
Dependencias & Se presentan mediante la columna de predecesoras y se explican como relación lógica entre actividades. \\
Duraciones & Se registran en la columna de duración para cada tarea y fase. \\
Recursos & Se incluyen en la hoja de recursos y en las asignaciones por actividad. \\
Hitos & Se identifican como tareas de duración cero, asociadas con aprobaciones y cierres parciales. \\
Ruta crítica & Se define como la secuencia de actividades que controla la fecha final del proyecto. \\
Responsables & Se indican por medio de los recursos asignados a cada tarea. \\
Costos & Se muestran mediante los costos fijos y el costo acumulado por tarea y fase. \\
\bottomrule
\end{tabularx}
\end{table}

\section{Configuración general del cronograma}

El cronograma se configuró con tareas programadas automáticamente y calendario estándar. A partir de la fecha de inicio definida, Microsoft Project calcula las fechas de comienzo y fin de las actividades según sus duraciones y predecesoras.

\begin{table}[H]
\centering
\caption{Configuración general del cronograma}
\begin{tabularx}{\textwidth}{p{5cm} Y}
\toprule
\textbf{Dato} & \textbf{Valor} \\
\midrule
Nombre del proyecto & Plataforma Universitaria Integrada \\
Herramienta utilizada & Microsoft Project \\
Fecha de inicio & 8 de junio de 2026 \\
Fecha final planificada & 9 de setiembre de 2026 \\
Duración aproximada & 3 meses \\
Presupuesto total & ₡17,625,200 \\
Método de programación & Tareas programadas automáticamente \\
Calendario & Estándar \\
\bottomrule
\end{tabularx}
\end{table}

\section{Evidencia visual del cronograma}

La evidencia visual del Diagrama de Gantt no se incrusta dentro de este documento para evitar pérdida de legibilidad o distorsión en las páginas horizontales. El cronograma visual exportado desde Microsoft Project se encuentra disponible en el repositorio del proyecto con el nombre \texttt{cronograma.pdf}. Dicho archivo permite observar las barras del calendario, las dependencias, la distribución temporal de las fases y la ruta crítica señalada por Microsoft Project.

De esta manera, este documento explica la estructura y los componentes del Entregable 6, mientras que el archivo \texttt{cronograma.pdf} funciona como evidencia visual directa del cronograma generado desde la herramienta.

\section{Estructura por fases}

El cronograma se organiza por fases con el fin de separar el trabajo de administración del proyecto, el desarrollo técnico del prototipo y las actividades de cierre. Esta estructura permite relacionar las tareas con los grupos de procesos estudiados en la materia.

\begin{table}[H]
\centering
\caption{Resumen de fases y costos del cronograma}
\begin{tabularx}{\textwidth}{p{3.8cm} Y p{3cm}}
\toprule
\textbf{Fase} & \textbf{Descripción} & \textbf{Costo} \\
\midrule
Inicio & Incluye aprobación del Business Case, elaboración del Project Charter y registro de stakeholders. & ₡1,886,667 \\
Planificación & Incluye alcance, WBS, cronograma, RACI, organigrama, recursos, costos, calidad, comunicaciones, riesgos y adquisiciones. & ₡3,686,667 \\
Ejecución & Incluye el desarrollo del prototipo Front-End navegable y sus portales principales. & ₡7,010,000 \\
Supervisión y control & Incluye KPIs, proceso de control de cambios, solicitudes de cambio, evaluación de cambios y dashboard ejecutivo del proyecto. & ₡2,271,667 \\
Cierre & Incluye informe ejecutivo, acta de cierre, reflexiones, recomendaciones y revisión final. & ₡588,333 \\
Rubros adicionales y contingencia & Incluye software, hardware, servicios operativos y reserva para contingencias. & ₡2,181,866 \\
\midrule
\textbf{Total} & Presupuesto total planificado del proyecto. & \textbf{₡17,625,200} \\
\bottomrule
\end{tabularx}
\end{table}

\section{Tareas, subtareas e hitos}

En Microsoft Project, las tareas resumen representan fases completas del trabajo. En este cronograma, las principales tareas resumen son Inicio, Planificación, Ejecución, Supervisión y control, Cierre, y Rubros adicionales y contingencia. Dentro de cada fase se incluyen tareas específicas con duración, fechas, predecesoras, responsables y costos.

Los hitos son actividades de duración cero que permiten marcar puntos importantes del proyecto sin agregar tiempo adicional al cronograma. Los hitos incluidos son los siguientes.

\begin{itemize}[leftmargin=*, label={\textbullet}]
    \item Hito: Project Charter aprobado.
    \item Hito: Alcance definido.
    \item Hito: Cronograma aprobado.
    \item Hito: Registro de riesgos completado.
    \item Hito: Prototipo Front-End completado.
    \item Hito: Cambios evaluados.
    \item Hito: Entrega final del proyecto.
\end{itemize}

\section{Dependencias y predecesoras}

Las predecesoras establecen la relación lógica entre actividades. Por ejemplo, el alcance debe quedar definido antes de iniciar tareas que dependen de esa definición, como el desarrollo del prototipo o la elaboración de ciertos documentos de planificación. Asimismo, el cierre del proyecto depende de la finalización del prototipo, la evaluación de cambios y la revisión final de los entregables.

Esta estructura permite que Microsoft Project recalcule automáticamente el cronograma si se modifica la duración o fecha de una tarea. Por ello, las dependencias ayudan a controlar el impacto que un atraso puede generar sobre actividades posteriores.

\section{Ruta crítica del proyecto}

La ruta crítica corresponde a la secuencia de actividades que determina la duración total del proyecto. Si una tarea de esta ruta se atrasa y no se recupera tiempo en otra parte, la fecha final del proyecto también puede atrasarse.

Con base en las dependencias registradas en Microsoft Project, la ruta crítica principal del proyecto se concentra en la definición del alcance, el desarrollo secuencial del prototipo Front-End, las pruebas de calidad y las actividades finales de cierre. La secuencia crítica queda conformada por las siguientes actividades.

\begin{itemize}[leftmargin=*, label={\textbullet}]
    \item Elaborar Project Charter.
    \item Hito: Project Charter aprobado.
    \item Definir alcance del proyecto.
    \item Hito: Alcance definido.
    \item Login y Dashboard principal.
    \item Portal Estudiantil - Matrícula y horarios.
    \item Portal Estudiantil - Calificaciones e historial.
    \item Portal Docente - Gestión de cursos y asistencia.
    \item Portal Docente - Registro de calificaciones.
    \item Portal Administrativo - CRUD completo.
    \item Integración entre módulos y navegación.
    \item Pruebas de usabilidad y QA.
    \item Hito: Prototipo Front-End completado.
    \item Elaborar informe ejecutivo.
    \item Elaborar acta de cierre y elaborar reflexiones y recomendaciones.
    \item Revisión final e integración de entrega.
    \item Hito: Entrega final del proyecto.
\end{itemize}

Aunque existen actividades importantes de planificación y supervisión que ocurren en paralelo, varias finalizan antes de las actividades principales de ejecución y cierre. Por esa razón, no todas las tareas del cronograma controlan directamente la fecha final.

\section{Recursos del proyecto}

La hoja de recursos contiene tanto recursos humanos como recursos de costo. Los recursos humanos representan a los integrantes que ejecutan o coordinan tareas, mientras que los recursos de costo representan rubros económicos que no se calculan por horas de trabajo.

\begin{table}[H]
\centering
\caption{Recursos humanos principales del proyecto}
\begin{tabularx}{\textwidth}{p{5cm} Y p{3.8cm}}
\toprule
\textbf{Recurso} & \textbf{Rol en el proyecto} & \textbf{Grupo} \\
\midrule
Tamara Robles Camacho & Project Manager & Gestión \\
Isaac Jiménez Blanco & Analista de Negocio y Documentador 1 & Documentación \\
Jorge Abarca Quiros & Analista de Negocio y Documentador 2 & Documentación \\
Mauricio González Prendas & Desarrollador Front-End 1 y QA & Desarrollo y QA \\
Carlos Sainz Arce & Desarrollador Front-End 2 y QA & Desarrollo y QA \\
\bottomrule
\end{tabularx}
\end{table}

\section{Costos y presupuesto}

El cronograma utiliza costos fijos por tarea. Esta decisión permite evitar una duplicación de costos por horas de recursos y mantener el presupuesto alineado con la estimación general del proyecto. Las tareas resumen muestran el costo acumulado de las actividades internas, lo que facilita el análisis de cuánto representa cada fase dentro del presupuesto total.

El presupuesto total planificado es de ₡17,625,200. La fase de ejecución concentra el mayor costo debido a que incluye el desarrollo del prototipo Front-End navegable. Además, se considera una reserva para contingencias que funciona como margen económico ante imprevistos menores.

\section{Relación con la materia}

El cronograma se relaciona directamente con contenidos vistos en el curso de administración de proyectos, ya que integra planificación del tiempo, asignación de recursos, estimación de costos, control de hitos y monitoreo de actividades críticas.

\begin{table}[H]
\centering
\caption{Relación del cronograma con los temas vistos en clase}
\begin{tabularx}{\textwidth}{p{4.2cm} Y}
\toprule
\textbf{Tema de la materia} & \textbf{Aplicación en el Entregable 6} \\
\midrule
WBS & La estructura de tareas se organiza por fases y paquetes de trabajo. \\
Cronograma & Se definen tareas, duraciones, fechas de inicio y fin. \\
Dependencias & Las predecesoras ordenan la secuencia lógica del proyecto. \\
Recursos & Se asignan responsables humanos y recursos de costo. \\
Costos & Se distribuye el presupuesto por actividad y fase. \\
Hitos & Se marcan puntos clave de aprobación o cierre. \\
Ruta crítica & Se identifica la secuencia que controla la fecha final. \\
Línea base & Permite comparar lo planificado contra lo ejecutado. \\
Monitoreo y control & El cronograma sirve para detectar atrasos, sobrecostos o desviaciones. \\
\bottomrule
\end{tabularx}
\end{table}

\section{Evidencias del entregable}

Para respaldar el Entregable 6 se deben conservar el archivo original de Microsoft Project y las evidencias exportadas desde la herramienta. Las evidencias principales son el archivo \texttt{.mpp}, la vista general del Diagrama de Gantt, la ruta crítica, la hoja de recursos, la vista de costos y las estadísticas del proyecto.

En el repositorio del proyecto, la evidencia visual principal del cronograma debe guardarse con el nombre \texttt{cronograma.pdf}, de forma que cualquier integrante pueda revisar la planificación completa sin depender únicamente del archivo editable de Microsoft Project.

\section{Conclusión}

El cronograma del proyecto Plataforma Universitaria Integrada permite visualizar la planificación completa de la iniciativa. Su valor principal consiste en organizar las actividades, asignar responsables, controlar costos, mostrar dependencias y definir la ruta crítica. Por esta razón, el Entregable 6 funciona como una herramienta de gestión y no solamente como un calendario de trabajo.

Asimismo, el cronograma sirve como base para el monitoreo y control del proyecto. Si durante la ejecución se presenta un atraso, una desviación de costo o un cambio importante, el equipo puede revisar su impacto dentro de Microsoft Project y tomar decisiones con mayor claridad. De esta manera, el cronograma contribuye a mantener el proyecto alineado con el alcance, el presupuesto, la calidad esperada y la fecha final planificada.


\newpage
\appendix

% =========================================================
% ANEXO A - TABLA DEL DIAGRAMA DE GANTT
% =========================================================

\begin{widepage}
\section{Tabla del Diagrama de Gantt}

La siguiente tabla corresponde a la vista del Diagrama de Gantt desde la columna Nombre de tarea hasta la columna Costo. Para mejorar la lectura, la información se divide por bloques de trabajo. En conjunto, estas tablas muestran duración, fechas, predecesoras, recursos asignados, costo fijo y costo total.

{\tiny
\renewcommand{\arraystretch}{1.08}
\begin{tabular}{L{4.2cm} L{1.35cm} L{1.75cm} L{1.75cm} L{1.35cm} L{6.1cm} L{1.55cm} L{1.55cm}}
\toprule
\textbf{Nombre de tarea} & \textbf{Duración} & \textbf{Comienzo} & \textbf{Fin} & \textbf{Predec.} & \textbf{Nombres de los recursos} & \textbf{Costo fijo} & \textbf{Costo} \\
\midrule
Inicio & 9 días & lun 6/8/26 & jue 6/18/26 & & & ₡0 & ₡1,886,667 \\
Aprobar Business Case & 5 días & lun 6/8/26 & vie 6/12/26 & & Isaac Jiménez Blanco, Tamara Robles Camacho & ₡850,000 & ₡850,000 \\
Elaborar Project Charter & 5 días & lun 6/8/26 & vie 6/12/26 & & Isaac Jiménez Blanco, Tamara Robles Camacho & ₡600,000 & ₡600,000 \\
Hito: Project Charter aprobado & 0 días & vie 6/12/26 & vie 6/12/26 & 3 & Tamara Robles Camacho & ₡0 & ₡0 \\
Registrar Stakeholders & 4 días & lun 6/15/26 & jue 6/18/26 & 2 & Isaac Jiménez Blanco, Tamara Robles Camacho & ₡436,667 & ₡436,667 \\
Planificación & 20 días & lun 6/15/26 & vie 7/10/26 & & & ₡0 & ₡3,686,667 \\
Definir alcance del proyecto & 5 días & lun 6/15/26 & vie 6/19/26 & 4 & Isaac Jiménez Blanco, Jorge Abarca Quiros & ₡350,000 & ₡350,000 \\
Hito: Alcance definido & 0 días & vie 6/19/26 & vie 6/19/26 & 7 & Isaac Jiménez Blanco, Jorge Abarca Quiros & ₡0 & ₡0 \\
Elaborar WBS & 5 días & lun 6/22/26 & vie 6/26/26 & 7 & Jorge Abarca Quiros, Isaac Jiménez Blanco & ₡450,000 & ₡450,000 \\
Crear cronograma del proyecto & 5 días & lun 6/29/26 & vie 7/3/26 & 9 & Isaac Jiménez Blanco, Jorge Abarca Quiros & ₡600,000 & ₡600,000 \\
Hito: Cronograma aprobado & 0 días & vie 7/3/26 & vie 7/3/26 & 10 & Isaac Jiménez Blanco & ₡0 & ₡0 \\
\bottomrule
\end{tabular}
}

\begin{center}
\textit{Tabla A.1. Vista tabular del Diagrama de Gantt: inicio y planificación base.}
\end{center}
\end{widepage}

\begin{widepage}
\section*{Continuación del Anexo A}

{\tiny
\renewcommand{\arraystretch}{1.08}
\begin{tabular}{L{4.2cm} L{1.35cm} L{1.75cm} L{1.75cm} L{1.35cm} L{6.1cm} L{1.55cm} L{1.55cm}}
\toprule
\textbf{Nombre de tarea} & \textbf{Duración} & \textbf{Comienzo} & \textbf{Fin} & \textbf{Predec.} & \textbf{Nombres de los recursos} & \textbf{Costo fijo} & \textbf{Costo} \\
\midrule
Elaborar matriz RACI & 3 días & lun 6/22/26 & mié 6/24/26 & 7 & Isaac Jiménez Blanco, Jorge Abarca Quiros & ₡150,000 & ₡150,000 \\
Crear organigrama del proyecto & 3 días & jue 6/25/26 & lun 6/29/26 & 12 & Isaac Jiménez Blanco, Jorge Abarca Quiros & ₡120,000 & ₡120,000 \\
Elaborar plan de recursos & 4 días & jue 6/25/26 & mar 6/30/26 & 12 & Isaac Jiménez Blanco, Jorge Abarca Quiros & ₡250,000 & ₡250,000 \\
Estimar costos del proyecto & 5 días & mié 7/1/26 & mar 7/7/26 & 14 & Isaac Jiménez Blanco, Jorge Abarca Quiros & ₡300,000 & ₡300,000 \\
Elaborar plan de calidad & 4 días & lun 6/22/26 & jue 6/25/26 & 7 & Isaac Jiménez Blanco, Jorge Abarca Quiros & ₡300,000 & ₡300,000 \\
Elaborar plan de comunicaciones & 4 días & mar 6/30/26 & vie 7/3/26 & 13 & Isaac Jiménez Blanco, Jorge Abarca Quiros & ₡250,000 & ₡250,000 \\
Elaborar plan de gestión de riesgos & 4 días & lun 6/22/26 & jue 6/25/26 & 7 & Isaac Jiménez Blanco, Jorge Abarca Quiros & ₡250,000 & ₡250,000 \\
Crear registro de riesgos & 4 días & vie 6/26/26 & mié 7/1/26 & 18 & Isaac Jiménez Blanco, Jorge Abarca Quiros & ₡250,000 & ₡250,000 \\
Hito: Registro de riesgos completado & 0 días & mié 7/1/26 & mié 7/1/26 & 19 & Isaac Jiménez Blanco & ₡0 & ₡0 \\
Crear matriz de probabilidad e impacto & 3 días & jue 7/2/26 & lun 7/6/26 & 19 & Isaac Jiménez Blanco, Jorge Abarca Quiros & ₡216,667 & ₡216,667 \\
Elaborar plan de adquisiciones & 3 días & mié 7/8/26 & vie 7/10/26 & 15 & Isaac Jiménez Blanco, Jorge Abarca Quiros & ₡200,000 & ₡200,000 \\
\bottomrule
\end{tabular}
}

\begin{center}
\textit{Tabla A.2. Vista tabular del Diagrama de Gantt: planificación complementaria.}
\end{center}
\end{widepage}

\begin{widepage}
\section*{Continuación del Anexo A}

{\tiny
\renewcommand{\arraystretch}{1.08}
\begin{tabular}{L{4.2cm} L{1.35cm} L{1.75cm} L{1.75cm} L{1.35cm} L{6.1cm} L{1.55cm} L{1.55cm}}
\toprule
\textbf{Nombre de tarea} & \textbf{Duración} & \textbf{Comienzo} & \textbf{Fin} & \textbf{Predec.} & \textbf{Nombres de los recursos} & \textbf{Costo fijo} & \textbf{Costo} \\
\midrule
Ejecución & 48 días & lun 6/22/26 & mié 8/26/26 & & & ₡0 & ₡7,010,000 \\
Ejecución - Desarrollo Front-End & 48 días & lun 6/22/26 & mié 8/26/26 & & & ₡0 & ₡7,010,000 \\
Login y Dashboard principal & 5 días & lun 6/22/26 & vie 6/26/26 & 8 & Carlos Sainz Arce, Mauricio González Prendas & ₡686,667 & ₡686,667 \\
Portal Estudiantil - Matrícula y horarios & 8 días & lun 6/29/26 & mié 7/8/26 & 25 & Mauricio González Prendas & ₡918,333 & ₡918,333 \\
Portal Estudiantil - Calificaciones e historial & 6 días & jue 7/9/26 & jue 7/16/26 & 26 & Mauricio González Prendas & ₡685,000 & ₡685,000 \\
Portal Docente - Gestión de cursos y asistencia & 6 días & vie 7/17/26 & vie 7/24/26 & 27 & Mauricio González Prendas & ₡710,000 & ₡710,000 \\
Portal Docente - Registro de calificaciones & 5 días & lun 7/27/26 & vie 7/31/26 & 28 & Mauricio González Prendas & ₡506,667 & ₡506,667 \\
Portal Administrativo - CRUD completo & 8 días & lun 8/3/26 & mié 8/12/26 & 29 & Mauricio González Prendas & ₡920,000 & ₡920,000 \\
Portal Financiero - Pagos e historial & 7 días & lun 6/29/26 & mar 7/7/26 & 25 & Carlos Sainz Arce & ₡706,667 & ₡706,667 \\
Dashboard Ejecutivo del sistema & 6 días & mié 7/8/26 & mié 7/15/26 & 31 & Carlos Sainz Arce & ₡666,667 & ₡666,667 \\
Integración entre módulos y navegación & 5 días & jue 8/13/26 & mié 8/19/26 & 30,32 & Carlos Sainz Arce, Mauricio González Prendas & ₡496,667 & ₡496,667 \\
Pruebas de usabilidad y QA & 5 días & jue 8/20/26 & mié 8/26/26 & 33 & Carlos Sainz Arce, Mauricio González Prendas & ₡713,332 & ₡713,332 \\
Hito: Prototipo Front-End completado & 0 días & mié 8/26/26 & mié 8/26/26 & 34 & Carlos Sainz Arce, Mauricio González Prendas & ₡0 & ₡0 \\
\bottomrule
\end{tabular}
}

\begin{center}
\textit{Tabla A.3. Vista tabular del Diagrama de Gantt: ejecución del prototipo.}
\end{center}
\end{widepage}

\begin{widepage}
\section*{Continuación del Anexo A}

{\tiny
\renewcommand{\arraystretch}{1.08}
\begin{tabular}{L{4.2cm} L{1.35cm} L{1.75cm} L{1.75cm} L{1.35cm} L{6.1cm} L{1.55cm} L{1.55cm}}
\toprule
\textbf{Nombre de tarea} & \textbf{Duración} & \textbf{Comienzo} & \textbf{Fin} & \textbf{Predec.} & \textbf{Nombres de los recursos} & \textbf{Costo fijo} & \textbf{Costo} \\
\midrule
Supervisión y control & 16 días & vie 6/26/26 & vie 7/17/26 & & & ₡0 & ₡2,271,667 \\
Definir KPIs del proyecto & 5 días & lun 7/6/26 & vie 7/10/26 & 10 & Jorge Abarca Quiros, Tamara Robles Camacho & ₡400,000 & ₡400,000 \\
Diseñar proceso de control de cambios & 4 días & vie 6/26/26 & mié 7/1/26 & 16 & Isaac Jiménez Blanco, Tamara Robles Camacho & ₡350,000 & ₡350,000 \\
Elaborar solicitud de cambio 1 & 3 días & jue 7/2/26 & lun 7/6/26 & 38 & Jorge Abarca Quiros, Tamara Robles Camacho & ₡250,000 & ₡250,000 \\
Elaborar solicitud de cambio 2 & 3 días & jue 7/2/26 & lun 7/6/26 & 38 & Jorge Abarca Quiros, Tamara Robles Camacho & ₡250,000 & ₡250,000 \\
Evaluar cambios solicitados & 4 días & mar 7/7/26 & vie 7/10/26 & 39,40 & Isaac Jiménez Blanco, Jorge Abarca Quiros & ₡350,000 & ₡350,000 \\
Hito: Cambios evaluados & 0 días & vie 7/10/26 & vie 7/10/26 & 41 & Tamara Robles Camacho & ₡0 & ₡0 \\
Crear dashboard ejecutivo del proyecto & 5 días & lun 7/13/26 & vie 7/17/26 & 37,42 & Jorge Abarca Quiros, Tamara Robles Camacho & ₡671,667 & ₡671,667 \\
Cierre & 10 días & jue 8/27/26 & mié 9/9/26 & & & ₡0 & ₡588,333 \\
Elaborar informe ejecutivo & 4 días & jue 8/27/26 & mar 9/1/26 & 35,42,43 & Jorge Abarca Quiros, Tamara Robles Camacho & ₡211,667 & ₡211,667 \\
Elaborar acta de cierre & 3 días & mié 9/2/26 & vie 9/4/26 & 45 & Isaac Jiménez Blanco, Tamara Robles Camacho & ₡150,000 & ₡150,000 \\
Elaborar reflexiones y recomendaciones & 3 días & mié 9/2/26 & vie 9/4/26 & 45 & Carlos Sainz Arce, Isaac Jiménez Blanco, Jorge Abarca Quiros, Mauricio González Prendas, Tamara Robles Camacho & ₡126,666 & ₡126,666 \\
Revisión final e integración de entrega & 3 días & lun 9/7/26 & mié 9/9/26 & 46,47 & Mauricio González Prendas, Tamara Robles Camacho, Carlos Sainz Arce, Isaac Jiménez Blanco, Jorge Abarca Quiros & ₡100,000 & ₡100,000 \\
Hito: Entrega final del proyecto & 0 días & mié 9/9/26 & mié 9/9/26 & 48 & Tamara Robles Camacho & ₡0 & ₡0 \\
\bottomrule
\end{tabular}
}

\begin{center}
\textit{Tabla A.4. Vista tabular del Diagrama de Gantt: supervisión y cierre.}
\end{center}
\end{widepage}

\begin{widepage}
\section*{Continuación del Anexo A}

{\tiny
\renewcommand{\arraystretch}{1.12}
\begin{tabular}{L{4.2cm} L{1.35cm} L{1.75cm} L{1.75cm} L{1.35cm} L{6.1cm} L{1.55cm} L{1.55cm}}
\toprule
\textbf{Nombre de tarea} & \textbf{Duración} & \textbf{Comienzo} & \textbf{Fin} & \textbf{Predec.} & \textbf{Nombres de los recursos} & \textbf{Costo fijo} & \textbf{Costo} \\
\midrule
Rubros adicionales y contingencia & 1 día & lun 6/8/26 & lun 6/8/26 & & & ₡0 & ₡2,181,866 \\
Software y herramientas & 1 día & lun 6/8/26 & lun 6/8/26 & & & ₡221,667 & ₡221,667 \\
Hardware y materiales & 1 día & lun 6/8/26 & lun 6/8/26 & & & ₡33,333 & ₡33,333 \\
Servicios operativos & 1 día & lun 6/8/26 & lun 6/8/26 & & & ₡411,000 & ₡411,000 \\
Reserva para contingencias & 1 día & lun 6/8/26 & lun 6/8/26 & & & ₡1,515,866 & ₡1,515,866 \\
\bottomrule
\end{tabular}
}

\begin{center}
\textit{Tabla A.5. Vista tabular del Diagrama de Gantt: rubros adicionales y contingencia.}
\end{center}
\end{widepage}

% =========================================================
% ANEXO B - HOJA DE RECURSOS
% =========================================================

\begin{widepage}
\section{Hoja de recursos}

La hoja de recursos muestra los recursos humanos, los recursos de costo y la forma en que se acumulan dentro del proyecto. Esta vista permite verificar que el cronograma incluye recursos de trabajo y recursos de costo, según lo requerido para el Entregable 6.

{\tiny
\renewcommand{\arraystretch}{1.18}
\begin{tabular}{L{4.6cm} L{1.8cm} L{1.5cm} L{1.5cm} L{3.2cm} L{1.5cm} L{1.9cm} L{1.9cm} L{1.6cm} L{1.8cm} L{2cm}}
\toprule
\textbf{Nombre del recurso} & \textbf{Tipo} & \textbf{Material} & \textbf{Iniciales} & \textbf{Grupo} & \textbf{Cap. máx.} & \textbf{Tasa estándar} & \textbf{Tasa h. extra} & \textbf{Costo/Uso} & \textbf{Acumular} & \textbf{Calendario base} \\
\midrule
Tamara Robles Camacho & Trabajo & & TRC & Gestión & 100\% & ₡0/hora & ₡0/hora & ₡0 & Prorrateo & Estándar \\
Isaac Jiménez Blanco & Trabajo & & IJB & Documentación & 100\% & ₡0/hora & ₡0/hora & ₡0 & Prorrateo & Estándar \\
Mauricio González Prendas & Trabajo & & MGP & Desarrollo y QA & 100\% & ₡0/hora & ₡0/hora & ₡0 & Prorrateo & Estándar \\
Jorge Abarca Quiros & Trabajo & & JAQ & Documentación & 100\% & ₡0/hora & ₡0/hora & ₡0 & Prorrateo & Estándar \\
Carlos Sainz Arce & Trabajo & & CSA & Desarrollo y QA & 100\% & ₡0/hora & ₡0/hora & ₡0 & Prorrateo & Estándar \\
Software y herramientas & Costo & & SW & Herramientas & & & & & Prorrateo & \\
Hosting / infraestructura & Costo & & HOST & Infraestructura & & & & & Prorrateo & \\
Hardware y materiales & Costo & & HW & Materiales & & & & & Prorrateo & \\
Servicios operativos & Costo & & SERV & Servicios & & & & & Prorrateo & \\
Contingencia & Costo & & CONT & Reserva & & & & & Prorrateo & \\
\bottomrule
\end{tabular}
}

\begin{center}
\textit{Tabla B.1. Hoja de recursos utilizada en Microsoft Project.}
\end{center}
\end{widepage}

\end{document}
